% Experimental Setup — auto-generated from benchmark outputs
\section{Experimental Setup}
\label{sec:experimental-setup}

We evaluate Semantic Fact DB against three baselines:
KnowledgeGraph (native triple store),
SheafDatabase (sheaf-theoretic representation),
and Apache Jena TDB2 (standard RDF store).

\subsection{Hardware}
All experiments run on a single machine with an
\texttt{\showcpu}
processor and
\texttt{\showmemory}
of RAM, running
\texttt{\showos}.

\subsection{Implementation}
Semantic Fact DB is implemented in Python 3.14.
The optimizer applies cost-based rewriting using
predicate pushdown, constant folding, and
dead-operator elimination.
Both the KG and Sheaf backends share the same
parser, logical plan representation, and optimizer.

\subsection{Datasets}
We use five datasets:
(1) Synthetic graphs with controlled size and arity,
(2) LUBM (Lehigh University Benchmark),
(3) DBpedia subset,
(4) YAGO subset,
(5) Wikidata subset.
Table~\ref{tab:datasets} summarizes their statistics.

% Dataset statistics — measured from benchmark generator
\begin{tabular}{@{}lrrr@{}}
\toprule
Dataset & Facts & Entities & Contexts \\
\midrule
Small    & 350   & 50  & 9 \\
Medium   & 1,630 & 500 & 51 \\
Large    & 17,150& 5,000 & 556 \\
\bottomrule
\end{tabular}


\subsection{Query Workloads}
Table~\ref{tab:workloads} lists the fifteen benchmark queries.
Queries Q1--Q4 test basic graph traversal,
Q5--Q7 test semantic filtering,
Q8 tests aggregation,
Q9--Q10 test mixed and global reconstruction,
Q11--Q13 test mutations,
and Q14--Q15 test I/O throughput.

% Workload descriptions (C1--C10)
\begin{tabular}{@{}llp{6cm}@{}}
\toprule
ID & Type & Description \\
\midrule
C1 & LOOKUP & Event reconstruction from open set membership \\
C2 & LOOKUP & Entity neighborhood traversal \\
C3 & LOOKUP & High-arity fact lookup (arity 3, 5, 8) \\
C4 & LOOKUP & High-arity multi-relational join \\
C5 & TEMPORAL & Temporal interval overlap query \\
C6 & TEMPORAL & Sliding-window aggregation \\
C7 & LOOKUP & Provenance chain traversal \\
C8 & GLOBAL & Consistency checking (cocycle condition) \\
C9 & GLOBAL & Global section construction (gluing) \\
C10 & MIXED & Multi-stage pipeline (future work) \\
\bottomrule
\end{tabular}


\subsection{Metrics}
For each query we measure:
latency (ms), memory (RSS bytes), CPU utilization,
and cache hit rate.
All metrics report mean, median, P95,
and 95\% confidence intervals over
10 (cold) or 50 (warm) runs.
